%%%%%%%%%%  scan idx 87  |  folio 74  %%%%%%%%%%
\textbf{LEMME 2.} \emph{Supposons qu'il existe sur \(X\) un fibré \(G\) à fibre vectorielle de rang \(p=\codim_XY\), tel que~:}
\begin{equation}\tag{3}
  \gamma_X(Y)=i_!(1)=\lambda_{-1}(G)\,;\qquad
  i^!(G)=\check E\,;\qquad
  (-1)^pC^p(G)=Y.
\end{equation}
\emph{Si alors on a~:}
\begin{equation}\tag{4}
  y=i^!(x),\qquad \text{\emph{avec} }x\in K(X),
\end{equation}
\emph{la formule (2) est valable.}

Démonstration immédiate, à partir du fait que \(\ch\) est un homomorphisme d'anneaux, et que l'on a~:
\begin{equation}\tag{5}
  \ch(\lambda_{-1}(G))=(-1)^pC^p(G)T(\check G)^{-1}.
\end{equation}

\textbf{COROLLAIRE.} \emph{\(\RR\) est vrai si \(Y\) est une hypersurface et si \(y\) est de la forme \(i^!(x)\), avec \(x\in K(X)\).}

\textbf{LEMME 3.} \emph{Supposons que \(X=Y\times P\), \(P\) étant un espace projectif, et que \(i\) soit l'injection \(u\longrightarrow (u,v_0)\) \(v_0\in P\). Alors \(i_*\) est injectif et \(\RR\) est vrai.}

C'est du pur calcul standard.

Supposons de nouveau \(X,Y,y\) quelconques (toujours avec \(Y\subset X\)). Soit \(X'\) la variété éclatée de \(X\) à l'aide de \(Y\), et \(Y'\) l'hypersurface image réciproque de \(Y\). On a un diagramme d'applications
\[
\begin{tikzcd}
Y \arrow[r,"i"] & X\\
Y' \arrow[u,"g"] \arrow[r,"j"] & X' \arrow[u,"f"']
\end{tikzcd}
\]
canoniques (\(j\) étant l'application canonique). L'image réciproque de \(E\) dans \(Y'\) admet un sous-fibré de rang \(1\), dont le dual est noté \(L\). D'ailleurs \(L^{-1}\) est aussi la restriction à \(Y'\) du fibré \(L(Y')\) défini sur \(X'\) par le diviseur \(Y'\). On pose \(F=E/L^{-1}\). Des calculs élémentaires donnent le~:

